% eksperymentalne tłumaczenie na włoski

%

\subsection{Symetrie środkowe} % lang-pl
\subsection{Simmetrie centrali} % lang-it

\begin{definition}
[symetria środkowa] % lang-pl
[simmetria centrale] % lang-it
    Niech $A = (x_A, y_A)$ będzie ustalonym punktem. % lang-pl
    Funkcję $S_A$ zadaną następująco: % lang-pl
    Sia $A = (x_A, y_A)$ un punto fissato. % lang-it
    Chiamiamo simmetria centrale di centro $A$ la funzione $S_A$ definita da: % lang-it
    \begin{equation}
        S_A(x_0, y_0) = (2x_A - x_0, 2y_A - y_0)
    \end{equation}
    nazywamy symetrią środkową o środku $A$. % lang-pl
\end{definition}

(Chodzi o to, by punkt $A$ był środkiem odcinka łączącego argument oraz wartość funkcji $S_A$). % lang-pl
(Vale a dire che il punto $A$ è il punto medio del segmento che unisce l'argomento e il valore della funzione $S_A$). % lang-it

\begin{example}
    Symetrie osiowe oraz środkowe to przykłady inwolucji. % lang-pl
    Le simmetrie assiali e le simmetrie centrali sono esempi di involuzioni. % lang-it
\end{example}

\begin{proposition}
    Złożenie dwóch symetrii środkowych jest translacją o wektor dwa razy dłuższy od tego, który łączy ich środki. % lang-pl
    La composizione di due simmetrie centrali è una traslazione di vettore doppio rispetto a quello che unisce i loro centri. % lang-it
\end{proposition}

Symetrie środkowe można emulować dwoma symetriami osiowymi: % lang-pl
Le simmetrie centrali possono essere ottenute mediante due simmetrie assiali: % lang-it

\begin{proposition}
\label{two_axial_one_point}
    Niech $k, l$ będą dwiema prostymi prostopadłymi do siebie, które przechodzą przez punkt $A$. % lang-pl
    Wtedy $S_k \circ S_l = S_l \circ S_k = S_A$. % lang-pl
    Siano $k$ e $l$ due rette perpendicolari tra loro passanti per il punto $A$. % lang-it
    Allora $S_k \circ S_l = S_l \circ S_k = S_A$. % lang-it
\end{proposition}

Równoważnie: jeżeli punkt $A$ leży na prostej $k$, to $S_A \circ S_k = S_k \circ S_A = S_l$, gdzie prosta $l$ jest prostopadła do prostej $k$ i przechodzi przez $A$. % lang-pl
Equivalentemente: se il punto $A$ appartiene alla retta $k$, allora $S_A \circ S_k = S_k \circ S_A = S_l$, dove la retta $l$ è perpendicolare a $k$ e passa per $A$. % lang-it
% Ćwiczenie 4.26, Neugebauer s. 203

\begin{proposition}
    Złożenie trzech symetrii środkowych jest symetrią środkową: $S_A \circ S_B \circ S_C = S_D$, gdzie punkt $D$ jest przesunięciem punktu $A$ o wektor $BC$. % lang-pl
    La composizione di tre simmetrie centrali è una simmetria centrale: $S_A \circ S_B \circ S_C = S_D$, dove il punto $D$ si ottiene traslando il punto $A$ del vettore $BC$. % lang-it
\end{proposition}

To jest, po krótkim zastanowieniu, treść twierdzenia Varignona. % lang-pl
\index{twierdzenie!Varignona}% % lang-pl
Mamy odpowiednik faktu \ref{klkm1}: % lang-pl
A ben vedere, questo è proprio il contenuto del teorema di Varignon. % lang-it
\index{teorema!di Varignon}% % lang-it
Abbiamo l'analogo del fatto \ref{klkm1}: % lang-it

\begin{proposition}
\label{klkm2}
    Niech $K, L$ będą punktami. % lang-pl
    Wtedy $S_K \circ S_L \circ S_K = S_M$, gdzie $M$ to punkt będący odbiciem punktu $L$ względem $K$. % lang-pl
    Siano $K$ e $L$ due punti. % lang-it
    Allora $S_K \circ S_L \circ S_K = S_M$, dove $M$ è il simmetrico del punto $L$ rispetto a $K$. % lang-it
\end{proposition}

Jeśli figura ma dwa środki symetrii, to ma ich nieskończenie wiele. % lang-pl
Se una figura possiede due centri di simmetria, allora ne possiede infinitamente molti. % lang-it

%